\chapter{Méthodes d’alignement}

« Apprendre dans un environement peu doté et l'une des directions clés de la recherche pour le TAL moderne. 
\citep{hedderich_2021}. » \citep{artemova_2023} Cette affirmation, présenté par \citet{artemova_2023} comme une
reformulation des travaux de \citet{hedderich_2021} viens confirmer ce que nous avons présenté auparavant dans les
chapitres de l'état de l'art. Maintenant la question n'est plus d'expliquer les problématiques et définitions
de notre environement mais plutôt les solutions élaborées par le passé. Ces méthodes d'alignement sont très 
pratiques et il se faut que nous les présentions par projet.

\section{Paracrawl et bitext}

ParaCrawl\footnote{https://paracrawl.eu/} \citep{banon_2020} est un robot d'exploration du web 
(et outil open-source) permettant le développement de larges corpus parallèles de textes traduit pour les 
langues officiels de l'Union Européenne. Des outils tel que Hunalign \citep{varga_2007} sont encore 
utilisés pour l'alignement de phrases pour des paires de document dans le projet ParaCrawl \citep{tan_2022}. 
Hunalign est un projet qui utilise des « heuristiques tels que la longueur des phrases 
ou la fréquence des mots pour trouver des paires de phrases dans deux documents. » \citep{tan_2022}

ParaCrawl mines des bitext\footnote{bitext} en passant par 4 étapes majeures : «
\begin{enumerate}
    \item Exploration de sites web : pour collecter du contenu provenant de pages web.
    \item Alignement documentaire : avec les pages collectées, trouvé le contenu qui s'aligne 
    dans différentes langues. Cette étape aligne le contenu au niveau documentaire.
    \item Alignement des phrases : avec les documents alignés, trouver des pair\-es de phrases correspondantes
    alignées.
    \item Filtrage des phrases : avec les phrases alignées, filtrer les paires de phrases de mauvaise qualité
    et utiliser le reste comme des données parallèles propre pour des taches en aval.
\end{enumerate} »
\citep{tan_2022}

Le projet de \citet{tan_2022} s'intéressent tout particuliérement aux étapes 3 et 4 et vise à les améliorer
par l'utilisation de meilleure fonction de notation. Ils essayent principalement de réaliser cette amélioration
pour une paire de langues peu dotées (ici Pashto et Khmer). Ainsi \citet{tan_2022} décident 
d'affiner le fonctionnement d'un modèles de transformateur de phrases afin d'aligner et de filtrer des phrases
pour les langues peu dotées recherchées. Pour ce faire ils utilisent un système d'apprentissage 
auto-supervisées \citep{chen_2020, henderson_2017} afin d'affiner le fonctionnement des modèles 
BERT \citep{devlin_2019} et \allowbreak SBERT \citep{reimers_2019}. Les résultats de leurs recherches montrent une
meilleures production de données d'entraînement (ici des bitextes) pour la traduction par réseau neuronal
qui aura lieu en aval de la récupération des données d'entraînement. Ces données ayant était récupérer
grâce au modèle SBERT \citep{reimers_2019} affiner avec Vecalign pour l'aligneur de phrases et un score basé 
sur la marge pour le filtrage des phrases alignées.

\section{BERT et SBERT}

BERT \citep{devlin_2019} (Bidirectional Encoder Representations from Transformers) et RoBERTa 
\citep{liu_2019} sont des modèles d'alignements de phrases préentrainé qui, à l'époque de leur créations, 
était très performant pour les tâches de régression de paires de phrases tel que la similarité 
sémantique\footnote{Similarité sémantique}. En France deux modèles se bases sur BERT : CamemBERT
\citep{martin_2020} et FlauBERT \citep{le_2020}. Ces deux modèles sont entrainé sur des données d'un 
corpus Français hétérogène, BERT ne se basant uniquement que sur des données d'entraînement en anglais.

Plus tard en 2019 \citet{reimers_2019} présente SBERT ou sentence-BERT qui se base sur le modèle préentrainé
BERT mais qui utilise des structures de réseau siamois et triplets\footnote{Réseau siamois}. Cette amélioration
réduit le temps pour trouver les paires de phrases les plus similaire à 5 secondes au lieu de 65 heures tout
en gardant la précision de BERT. Après évaluation SBERT surpasse les autres méthodes de plongement de phrases
\footnote{plongement de phrases (sentence embedding)} pour des tâches de similarité sémantique communes.
\citep{reimers_2019}

\section{CLIR et plongement lexical}

FaDA (fast document aligner - aligneur de document rapide) \citep{lohar_2016} est un outil d'alignement de 
documents multilingues se basant sur un alogrithme d'alignement documentaire CLIR\footnote{CLIR} 
(crosslingual information retrieval - récupération d'information interlangue). « Un système d'aligneur de 
document interlangue vise à extraire efficacement des candidats probables de documents alignés venant d'un
corpus comparable dans deux ou plus langues. » \citep{lohar_2016} D'autres travaux existent dans l'état
de l'art tel que le projet \citet{zhao_2002} permettant le minage de phrases parallèles de corpus 
comparables bilingues. Ou le projet STRAND par \citet{resnik_2003} présentant un système basé sur du
minage sur des sites web. D'autres projets présentes des méthodes se basant sur différents niveaux comme
les travaux de \citet{yang_2003} ou encore le système de requète sur bases de données pour des lexiques
bilingues présenté par \citet{munteanu_2005}.

\citet{lohar_2016} propose une méthode de plongement lexical\footnote{plongement lexical} vectoriel 
incorporant une approche basé sur la CLIR pout estimer la similarité sémantique entre deux documents de
langue source et langue cible. Leur système commence par indexer les documents de langue source et langue
cible. Les documents sources sont utilisé pour faire une pseudo-requête composé d'une fraction des termes
constituant les documents dans la langue source. Les termes choisis sont trouvé en se basant sur un score
se basant sur une équation utilisant la fréquence des termes, le nombre total de documents et le nombre
de documents dans lesquels le terme apparaît. Les termes en langue source sont ensuite traduit via un 
dictionnaire vers la langue cible indexé avant d'être utiliser pour réaliser une comparaison entre la liste
de terme et les documents dans la langue cible. Cette méthode de comparaison permet d'extraire les meilleures
documents avec de sélectionner celui avec le meilleur score pour avoir le meilleur candidat pour l'alignement.

\citet{steingrimsson_2021} utilise FaDA ainsi qu'une méthode de notation de phrases candidates basé sur
le plongement contextualisé et d'autre alignement par mot de haute précision. Enfin un classificateur 
sélectionne des paires de phrases en se basant sur ces scores.

Il est intéressant de noter que le papier de \citet{steingrimsson_2021} se base sur des données venant d'ensemble
de tests compilés ressemblant à ceux réaliser pendant l'atelier BUCC 2017 \citep{zweigenbaum_2016} et sur
des un test basé sur des articles wikipedia, dans deux langues étudiés puis décomposés en phrases, et les 
outils CLIR utilisés pour obtenir les traductions candidates.

\section{Réutilisation de segment parallèles}

\citet{steingrimsson_2023} ont réalisé des travaux de recherches sur la réexamination de données d'entraînement
parallèles rejetés en les segmentant avant d'extraire sementiquement d'autres segments bilingues equivalents.
Ces segments sont gadés si ceux-ci sont reconnu comme utile pour l'entraînement de traducteur automatique selon
la méthode présentés dans leurs travaux. La qualité de la sortie est comparé à celle d'un modèle référentiel.
En utilisant cette technique le temps d'entraînement est réduit de 65\% tout en augmentant le score BLEU de
1,6. Ce système, entrainé sur moins de données montre de meilleures résultat que le référentiel selon des
tests fait avec MultEval \citep{clark_2011}.